% !TeX root = ../../Book1.tex
% 纲要标题：### 17.4 有限域的计算与编码
% 纲要位置：计划纲要.md，第 606 行。

\section{有限域的计算与编码}
\label{b1:sec:1704}

% 纲要内容（仅作写作计划，尚非教材正文）：
%
% * 最小多项式、轨道和与轨道积；区分不同共轭元与全部嵌入像的计数
% * 多项式分解的有限域方法；不可约性检验与重复因子的重数恢复
% * Hamming 距离、求值码与唯一纠错；进一步的编码算法置于附录
%
% ---
%

Frobenius 轨道列出有限域元素的共轭元，由此可以计算最小多项式，并解释共轭元的和与积为何回到基域。下一章将从乘法线性算子的迹与行列式定义域迹、域范数，再证明它们在有限域扩张中可以用全部嵌入像的和与积计算，本节先计算这些和与积。有限域运算还可以构造纠错码：多项式在互异点处取值，根数界便控制不同编码向量之间的差别。

\subsection{最小多项式、迹与范数的计算}
\label{b1:subsec:1704:01}

\begin{proposition}[Frobenius 轨道与最小多项式]\label{b1:prop:frobenius-orbit-minimal}
设 \(q\) 为素数幂、\(n\geq1\) 为整数，\(\alpha\in\mathbb F_{q^n}\)。令 \(d\) 为使 \(\alpha^{q^d}=\alpha\) 的最小正整数。则 \(d\mid n\)，而 \(\alpha\) 在 \(\mathbb F_q\) 上的首一最小多项式为
\begin{equation}\label{b1:eq:frobenius-minimal-polynomial}
m_{\alpha,\mathbb F_q}(x)=\prod_{j=0}^{d-1}\bigl(x-\alpha^{q^j}\bigr).
\end{equation}
特别地，\([\mathbb F_q(\alpha):\mathbb F_q]=d\)，即元素的扩张次数等于其 Frobenius 轨道长度。
\end{proposition}
\begin{proof}
令 \(F:\mathbb F_{q^n}\rightarrow\mathbb F_{q^n}\) 为 \(q\) 次 Frobenius 自同构，\(F(z)=z^q\)。它满足 \(F^n=1\)，不动元恰为 \(\mathbb F_q\)。将 \(n\) 除以 \(d\)，写成 \(n=kd+r\)、\(0\leq r<d\)。由于 \(F^d(\alpha)=\alpha\)，有
\[
\alpha=F^n(\alpha)=F^r(\alpha).
\]
若 \(r>0\)，这与 \(d\) 是最小正周期矛盾，故 \(r=0\)，即 \(d\mid n\)。轨道中的 \(\alpha,\alpha^q,\ldots,\alpha^{q^{d-1}}\) 两两不同。

记右端乘积为 \(P(x)\)。对 \(P\) 的各个系数施行 \(F\)，并保持形式变量 \(x\) 不变，得到
\[
F(P)(x)=\prod_{j=0}^{d-1}\bigl(x-\alpha^{q^{j+1}}\bigr)=P(x),
\]
因为 \(F\) 只是把这 \(d\) 个根循环置换。因此每个系数 \(c\) 都满足 \(c^q=c\)，所以 \(P\in\mathbb F_q[x]\)。

另一方面，若非零 \(h\in\mathbb F_q[x]\) 满足 \(h(\alpha)=0\)，其系数都被 \(F\) 固定，故
\[
h(\alpha^q)=h(\alpha)^q=0.
\]
迭代此等式，轨道中的全部 \(d\) 个元素都是 \(h\) 的根，因而 \(\deg h\geq d\)。\(P\) 首一、以 \(\alpha\) 为根且次数为 \(d\)，于是它就是最小多项式；\cref{b1:prop:minimal-polynomial-degree}再给出所述扩张次数。
\end{proof}

在 \(\mathbb F_8=\mathbb F_2(\alpha)\)、\(\alpha^3=\alpha+1\) 中，元素在 \(\mathbb F_2\) 上的次数只能是三的除数。除 \(0,1\) 外，剩下六个元素的次数都为三，故平方映射将它们分成两个长度为三的轨道。全部轨道如下：
\begin{center}
\begin{tabular}{@{}lcl@{}}
\toprule
轨道 & 长度 & 对应的首一最小多项式 \\
\midrule
\(\{0\}\) & \(1\) & \(x\) \\
\(\{1\}\) & \(1\) & \(x+1\) \\
\(\{\alpha,\alpha^2,\alpha^4\}\) & \(3\) & \(x^3+x+1\) \\
\(\{\alpha^3,\alpha^6,\alpha^5\}\) & \(3\) & \(x^3+x^2+1\) \\
\bottomrule
\end{tabular}
\end{center}
由于 \(\alpha^3=\alpha+1\)，可核对 \(\alpha^3\) 满足 \(x^3+x^2+1\)；两个三次式在 \(\mathbb F_2\) 中都无根，故不可约。四个轨道互不相交，长度之和为 \(1+1+3+3=8\)，每个轨道对应一个不可约因子。因此
\[
x^8-x=x(x+1)(x^3+x+1)(x^3+x^2+1),
\]
并且 \(N_2(3)=2\)。\(\alpha\) 的三个轨道元素之和为零，积为 \(\alpha^7=1\)；在 \(\mathbb F_8/\mathbb F_2\) 中它们也是 \(\alpha\) 的全部嵌入像，下一章的迹与范数公式将分别给出零和一。

\ProseParagraphBreak
元素所处的扩张不能省略。例如把 \(1\) 视为 \(\mathbb F_4/\mathbb F_2\) 的元素，按扩张的两个嵌入计，和为 \(1+1=0\)、积为一；只按它的单元素轨道计则和为一。一般地，设 \(q\) 为素数幂、\(n\geq1\)，记 \(F(z)=z^q\)，若 \(\alpha\in\mathbb F_{q^n}\) 的轨道长为 \(d\)，则 \(d\mid n\)。该扩张的各个嵌入的像都是代数闭包中同一个 \(q^n\) 元子域，所以它们正是\cref{b1:thm:finite-field-automorphisms}列出的 \(n\) 个自同构 \(F^0,\ldots,F^{n-1}\)。这些嵌入在 \(\alpha\) 上的取值每经过 \(d\) 步便重复一次，所以每个轨道元素出现 \(n/d\) 次。下一章的迹公式将轨道元素的和重复 \(n/d\) 次相加，范数公式则将轨道元素的积取 \(n/d\) 次幂，不能总以最小多项式的一组不同根的和与积代替。

\subsection{有限域上的多项式分解}
\label{b1:subsec:1704:02}

设 \(q\) 为素数幂。非零 \(f\in\mathbb F_q[x]\) 若每个首一不可约因子都只出现一次，就称为无重因子多项式。\cref{b1:prop:perfect-frobenius}已说明有限域是完美域，即有限域上的不可约多项式都可分，所以对非常数 \(f\)，这等价于 \(f\) 在代数闭包中没有重根。由\cref{b1:prop:finite-field-factorization}，对这样的 \(f\) 计算
\[
\gcd\bigl(f,x^{q^d}-x\bigr)
\]
就能提取 \(f\) 中次数整除 \(d\) 的全部不可约因子。例如 \(d=6\) 时，这个最大公因式可能同时含有一次、二次、三次和六次因子，不能直接当作六次因子的乘积。幂 \(x^{q^d}\) 可以不断取 \(q\) 次幂并模 \(f\) 取余，无须展开到巨大次数。按 \(d\) 从小到大处理，从本次最大公因式中除去其中已经找出的低次因子，便得到次数恰为 \(d\) 的因子乘积。

\begin{proposition}[有限域上的不可约性检验]\label{b1:prop:finite-field-irreducibility-test}
设 \(q\) 为素数幂，\(f\in\mathbb F_q[x]\) 为次数 \(n>1\) 的首一多项式。则 \(f\) 不可约，当且仅当 \(f\mid x^{q^n}-x\)，并且对每个整除 \(n\) 的素数 \(\ell\)，都有
\[
\gcd\bigl(f,x^{q^{n/\ell}}-x\bigr)=1.
\]
\end{proposition}
\begin{proof}
若 \(f\) 不可约，\cref{b1:prop:finite-field-factorization}给出 \(f\mid x^{q^n}-x\)。对整除 \(n\) 的素数 \(\ell\)，有 \(n/\ell<n\)，故 \(n\) 不整除 \(n/\ell\)，于是 \(f\) 不整除 \(x^{q^{n/\ell}}-x\)。由 \(f\) 不可约，相应最大公因式只能为一。

反过来，\(x^{q^n}-x\) 的导数为 \(-1\)，所以条件 \(f\mid x^{q^n}-x\) 首先排除了 \(f\) 的重复因子。\cref{b1:prop:finite-field-factorization}还说明 \(f\) 的每个不可约因子 \(h\) 的次数 \(d\) 都整除 \(n\)。若 \(d<n\)，整数 \(n/d>1\) 有某个素因子 \(\ell\)，从而 \(n/d=\ell s\)，即 \(n/\ell=ds\)。因此 \(d\mid n/\ell\)，同一命题给出 \(h\mid x^{q^{n/\ell}}-x\)，与该最大公因式为一矛盾。这说明每个不可约因子的次数都等于 \(n\)；而 \(\deg f=n\)，故只有一个因子，\(f\) 不可约。
\end{proof}

处理任意非零 \(f\in\mathbb F_q[x]\) 时，先提出首项系数，对剩余的首一多项式递归；遇到常数一便停止。若 \(f'\ne0\)，记 \(h=\gcd(f,f')\)。当 \(h=1\) 时，\(f\) 已无重因子；当 \(h\ne1\) 时，须分别递归处理 \(h\) 和 \(f/h\)，二者次数都严格小于 \(\deg f\)。不能只处理商 \(f/h\)，因为它可能遗漏只在 \(h\) 中留下的不可约因子。若 \(f'=0\)，则所有非零项的指数都被特征 \(p\) 整除，可写成 \(f=\sum_i a_ix^{pi}\)。由\cref{b1:prop:perfect-frobenius}，有限域上的 Frobenius 满射，并且仍为单射，所以每个系数 \(a_i\) 有唯一的 \(p\) 次根 \(b_i\)。于是
\[
g=\sum_i b_ix^i\qquad\text{满足}\qquad f=g^p.
\]
递归处理 \(g\) 后，将所得因子的重数都乘以 \(p\)。不同分支得到相同首一不可约因子时，还须把重数相加。每次递归都降低正次数，因此过程终止；每一步保持乘积等式，所以因子及其重数均被保留。\subsecref{b1:subsec:B01:02}给出这一预处理的进一步算例。

\ProseParagraphBreak
例如在 \(\mathbb F_2[x]\) 中取 \(g=x\)、\(h=x+1\)，令 \(f=g^3h^2\)。特征为二，故
\[
f'=g^2h^2,\qquad \gcd(f,f')=g^2h^2,\qquad
\frac{f}{\gcd(f,f')}=g.
\]
商中完全没有 \(h\)，若只分解它就会漏掉这个因子。另一分支 \(g^2h^2=(gh)^2\) 的导数为零，开平方得到无重因子的 \(gh\)，再把两个一次因子的重数乘以二。与商分支的一个 \(g\) 合并后，\(g\) 的重数为 \(2+1=3\)，\(h\) 的重数为二，恢复原来的 \(g^3h^2\)。

\ProseParagraphBreak
按次数分组不一定已经得到每个不可约因子。对于同次数因子的乘积，可以枚举该次数的首一多项式并作试除；枚举有限且必终止。这说明原则上能通过有限搜索完成分解，但候选数随次数增长很快，实际计算还需要改善效率。\subsecref{b1:subsec:B01:03}进一步给出 Berlekamp 分解\footnote{伯利坎普（Elwyn Ralph Berlekamp，1940-2019），美国数学家，研究代数编码与组合博弈，提出有限域多项式的分解算法。}，用线性代数和最大公因式分开这些因子。

\subsection{编码理论的应用}
\label{b1:subsec:1704:03}

\begin{definition}\label{b1:def:hamming-distance}
设 \(q\) 为素数幂，\(N\geq1\) 为整数。对 \(v,w\in\mathbb F_q^N\)，数
\[
d_H(v,w)\coloneqq\#\{\,i\in\{1,\ldots,N\}\mid v_i\ne w_i\,\}
\]
称为 \(v,w\) 的\term[Hamming-juli]{Hamming 距离}[Hamming distance]。\footnote{汉明（Richard Wesley Hamming，1915-1998），美国数学家与计算机科学家，建立纠错码理论中的汉明码及相关距离概念。}若 \(C\subseteq\mathbb F_q^N\) 至少包含两个向量，则不同向量之间 Hamming 距离的最小值称为 \(C\) 的\term[zuixiao-juli]{最小距离}[minimum distance]，记为 \(d(C)\)。
\end{definition}
\symidx{\(d_H(v,w)\)：Hamming 距离}
\symidx{\(d(C)\)：码的最小距离}
距离为零恰好表示两个向量相同，距离也显然对称。若 \(v_i\ne w_i\)，则对任意向量 \(z\)，必有 \(v_i\ne z_i\) 或 \(z_i\ne w_i\)。对这些坐标计数，便有三角不等式
\[
d_H(v,w)\leq d_H(v,z)+d_H(z,w).
\]

\begin{samepage}
下面构造一个\term[qiuzhi-ma]{求值码}[evaluation code]。取 \(\mathbb F_5\) 中五个互异点 \(0,1,2,3,4\)。对信息 \((a,b)\in\mathbb F_5^2\)，取 \(f(x)=a+bx\in\mathbb F_5[x]\)；它可以是零多项式，否则次数不超过一。编码向量为
\[
\bigl(f(0),f(1),f(2),f(3),f(4)\bigr)\in\mathbb F_5^5.
\]
\end{samepage}
记全部编码向量组成的集合为 \(C\)，其中的向量称为码字。若两个信息对应不同的多项式 \(f,g\)，则 \(h=f-g\ne0\)，且 \(\deg h\leq1\)。在五个互异求值点中，至多一个点满足 \(h(a)=0\)，也就是 \(f(a)=g(a)\)。所以相应码字至少在四个坐标不同，即 \(d(C)\geq4\)。零多项式与 \(x\) 的取值只在点零处相同，给出距离恰为四的一对码字，因此 \(d(C)=4\)。若传输中至多改坏一个坐标，原码字便唯一：若两个码字 \(c,c'\) 都与收到的向量 \(y\) 相距至多一，三角不等式便给出 \(d_H(c,c')\leq2<4\)，与最小距离矛盾。

\ProseParagraphBreak
例如 \(f(x)=1+2x\) 给出 \((1,3,0,2,4)\)。若收到 \((1,3,1,2,4)\)，可枚举所有形如 \(a+bx\) 的多项式（包括零多项式），找出在至少四个位置吻合者，唯一得到原来的 \(f\)。也可以枚举五个求值点的所有两点子集，每次由相应的两个接收值插值得到一个次数不超过一的多项式，再检验它与全部五个接收值的吻合数。由于至多一处错误，必有选中的两个位置都正确，此时插值恢复原来的 \(f\)；包含错误位置的候选也须通过至少四处吻合的检验，不能任取两个位置就相信所得结果。前一种方法只有二十五个候选，后一种方法只需尝试十组点，两者都给出完整的有限解码步骤；更长的码需要更有效的方法。这一例子说明了多项式根数界如何转化为纠错能力。
